\section*{PROJECT SUMMARY}
\phantomsection\addcontentsline{toc}{section}{PROJECT SUMMARY}

This thesis designs and implements an FPGA accelerator for the inference stage of a quantized HiFiC GAN, a Generative Adversarial Network used for learned image compression. The work follows a hardware--software co-design methodology: the computationally dominant residual blocks of the HiFiC generator are accelerated on the Programmable Logic, while model control and auxiliary functions run in software on the Processing System.

The main contributions of the thesis are summarized as follows:
\begin{itemize}
    \item A post-training mixed-precision quantization of the HiFiC model, with FP32 input/output interfaces and FP16 internal computation, that preserves compression quality: PSNR, SSIM, and LPIPS remain almost unchanged compared with the FP32 baseline.
    \item A full-generator accelerator built with High-Level Synthesis, organized into a Fusion core (initial convolution, nine residual blocks, and the global skip), an up-convolution core, and a final $7\times7$ convolution core, using line buffers, parallel processing-element clusters, and a rotating accumulator to maximize data reuse and on-chip parallelism.
    \item A complete implementation on the ZCU104 platform operating at 300~MHz (maximum 308.74~MHz) that reconstructs one $256\times256$ image in 735.6~ms at a throughput of 122.8~GOPS, about $1459\times$ faster than a single-threaded CPU software reference.
\end{itemize}

These results show that FPGA-based acceleration is a practical and effective approach for deploying GAN-based learned image compression, and provide a foundation for more comprehensive hardware deployment in future work.

\cleardoublepage
